% Section I -- Introduction.
%
% NOVELTY CLAIMS: the phrases "to our knowledge" below need checking against
% the literature before submission. Verify specifically whether any prior work
% measures the energy of DP-FL rather than modelling it, and whether the
% per-round / convergence decomposition has been made elsewhere. Soften or
% drop the claims that do not survive that check -- a reviewer who knows one
% counterexample will discount the rest of the paper.
% Citation keys are placeholders: \cite{dpsgd}, \cite{fedavg}, \cite{opacus},
% \cite{flower}, \cite{greenai}, \cite{strubell}, \cite{rdp}, \cite{uci-har}.

\section{Introduction}

Differential privacy in federated learning is evaluated almost exclusively
along two axes. A privacy budget $\varepsilon$ is chosen, the resulting
accuracy is reported, and the trade-off between them is the result. The
literature on DP-SGD~\cite{dpsgd} and its federated variants has made that
trade-off precise, and tooling such as Opacus~\cite{opacus} has made it
routine to measure.

A third axis is missing. Training consumes energy, and adding differential
privacy changes how much. On the battery-powered edge devices that federated
learning exists to serve, energy is not a secondary consideration to be
optimised after the fact --- it decides whether a deployment is possible at
all. A privacy configuration that a phone cannot afford to run is not a point
on the privacy--utility curve; it is off the curve entirely.

Green AI research has established that the energy and carbon costs of machine
learning deserve first-class measurement~\cite{greenai, strubell}, but that
work has concentrated on large centralised training. The intersection ---
what privacy costs in joules, in a federated setting --- has not, to our
knowledge, been measured directly.

This paper measures it, and reports three findings that a two-axis evaluation
cannot produce.

\textbf{The dominant cost of privacy is convergence, not arithmetic.} Adding
differential privacy imposes a fixed per-round overhead from per-sample
gradient computation and clipping: we measure $+45\%$ on UCI-HAR and $+60\%$
on CIFAR-10, and this overhead does not vary with $\varepsilon$ within our
measurement resolution. It is the visible cost, and it is the smaller one.
Everything that depends on the privacy budget acts instead through the number
of rounds required, which an energy-per-round evaluation cannot see. The two
mechanisms have opposite engineering implications: a fixed overhead is
addressed by optimising the DP implementation, a convergence penalty by
changing the round budget.

\textbf{The cost grows with the accuracy demanded, and becomes unbounded.} On
HAR the private and non-private conditions cost the same energy to reach
$0.40$ accuracy, $1.2$--$2.1\times$ to reach $0.50$, and $5.8\times$ to reach
$0.70$, with the tighter budgets never arriving. At $0.85$ accuracy, which the
non-private model reaches comfortably, no privacy budget we tested arrives at
any energy expenditure. A privacy budget therefore implies a ceiling on
achievable model quality, past which additional computation buys nothing ---
because under client-level accounting the additional training is itself what
raises the noise.

\textbf{The same privacy budget carries very different carbon costs.} Because
the cost is paid in joules, it inherits the carbon intensity of the grid
supplying them. The privacy premium on 1000 HAR training runs is
$105$--$133$~gCO$_2$eq on Algeria's gas-dominated grid against $5.6$--$7.1$~g
in Norway. Selecting $\varepsilon$ is in part a geographic decision, which the
privacy--utility framing has no way to express.

We measure energy with NVML hardware counters differenced around each
federated round, rather than estimating it from device power ratings as much
Green AI work does. Two datasets are used. UCI-HAR~\cite{uci-har} is primary:
its clients are the individual volunteers the data was collected from, so the
federation's client boundary is a real person rather than a synthetic
partition --- an arrangement that matters for a study about protecting
individuals' data, and one that places the work in the sensor-on-edge setting
where these deployments live. CIFAR-10 is secondary and checks the same
effects on a harder benchmark.

A methodological contribution follows from the measurement itself. Deriving
the noise multiplier requires a search over the privacy accountant, and that
search costs more for a loose budget than a tight one. Executed inside the
federated loop, where the natural implementation puts it, its cost lands
inside the measurement window and scales with $\varepsilon$ --- taking exactly
the shape of a cost of privacy. In our first clean sweep it accounted for the
entire observed trend. We describe this artefact and the three controls that
remove it, because any energy study of a privacy mechanism is exposed to it.

The rest of the paper is organised as follows. Section~II reviews related
work. Section~III describes the accounting model, the energy measurement and
the experimental controls. Section~IV reports results on both datasets.
Section~V discusses the implications, and Section~VI concludes.
